\documentclass{article}
\usepackage{amsmath,amssymb}
\usepackage[margin=1in]{geometry}
\allowdisplaybreaks
\usepackage[UTF8]{ctex}
\begin{document}

\section*{一般选择型有限视界收益管理动态规划}

\begin{equation}
V_t(x)=\max_{k=1,\ldots,m}
\left\{\lambda\bigl(R_k-Q_k\Delta V_{t-1}(x)\bigr)\right\}
+V_{t-1}(x).
\tag{7}
\end{equation}
递归过程在状态 $t=1,\ldots,T$ 和 $x=1,\ldots,C$ 处发生，同时也会继承原有的边界条件。
\begin{equation}
V_t(0)=0,\quad t=1,\ldots,T
\qquad\text{and}\qquad
V_0(x)=0,\quad x=1,\ldots,C.
\tag{2}
\end{equation}
在每个这样的状态下，做出的决策是一个索引$k\in\{1,\ldots,m\}$：所提供的集合是$S_k$。在(7)中，一个最大化索引可以提供一个最优的提供集合。这些集合、它们的排序以及它们的系数在下面有定义；它们并非任意输入的。

\paragraph{状态、固定数据以及解释（参见~第5-7页）。}
有$n$种票价产品，这些产品的票价都是固定的。
\[
N=\{1,\ldots,n\},
\qquad r_1\ge r_2\ge\cdots\ge r_n\ge0.
\]
$C$表示库存中的飞机数量，而$T$则表示时间周期的数量。状态$x$表示剩余的库存数量，而$t$则表示剩余的时间周期数量。时间是以离散的方式计算的，且是从后往前计算的：$t=0$表示销售截止日期。因此，状态$x=0,\ldots,C$和$t=0,\ldots,T$分别表示不同的时间范围。数值$V_t(x)$表示在周期$t,t-1,\ldots,1$期间，考虑到在时间$t$时还有$x$台飞机剩余的情况下，可以获得的最大预期收入。

每个时间段内，最多只有一位顾客到达。固定的到达概率$\lambda$在每个时间段都是相同的。该公司提供了一组产品$S\subseteq N$。在顾客到达的条件下，$P_j(S)$是选择产品$j\in S$的概率；$j=0$表示不购买，而$P_0(S)$则是购买该产品的概率。这些固定的选择概率不依赖于时间。选择结果仅取决于当前提供的产品组合，而不受卖方政策、其他买家的策略或购买历史的影响。作者要求有一个合适的概率函数来描述这种情况。
\begin{align*}
P_j(S)&=0,
&&j\in N\setminus S,\quad S\subseteq N,\\
P_j(S)&\ge0,
&&j\in S,\quad S\subseteq N,\\
\sum_{j\in S}P_j(S)+P_0(S)&=1,
&&S\subseteq N.
\end{align*}
对于所提供的集合$S$，周期概率如下：
\begin{align*}
\Pr(\text{sale of product }j\text{ in the period}\mid S)
&=\lambda P_j(S),\qquad j\in N,\\
\Pr(\text{no sale in the period}\mid S)
&=\lambda P_0(S)+(1-\lambda).
\end{align*}
第二个表达式中包含了“无到达”和“有到达但无购买”两种情况。

\paragraph{边际值和综合系数（p.~7）。}
在(7)中使用的容量边际成本为
\[
\Delta V_{t-1}(x)=V_{t-1}(x)-V_{t-1}(x-1),
\qquad t=1,\ldots,T,\quad x=1,\ldots,C.
\]
对于每一个提供的集合$S\subseteq N$，进行如下定义：
\begin{align*}
Q(S)&=\sum_{j\in S}P_j(S)=1-P_0(S),\\
R(S)&=\sum_{j\in S}P_j(S)r_j.
\end{align*}
这里的$Q(S)$表示购买的总概率，而$R(S)$则是指在某种条件下，提供$S$所能带来的总预期收入。到达因子则被保留在公式(7)中的系数之外。

\paragraph{非支配家族的构建：定义 1（p.~8）。}
仅在这一定义中，$T$表示$N$的一个候选子集，而非地平线长度。如果存在满足某些子集的概率$\alpha(S)$，则称集合$T\subseteq N$为被支配的。
\[
\alpha(S)\ge0\quad(S\subseteq N),
\qquad \sum_{S\subseteq N}\alpha(S)=1,
\]
使得以下情况之一发生：
\[
Q(T)\ge\sum_{S\subseteq N}\alpha(S)Q(S)
\quad\text{and}\quad
R(T)<\sum_{S\subseteq N}\alpha(S)R(S),
\]
或者
\[
Q(T)>\sum_{S\subseteq N}\alpha(S)Q(S)
\quad\text{and}\quad
R(T)\le\sum_{S\subseteq N}\alpha(S)R(S).
\]
如果一个集合没有被其他集合所支配，那么就说这个集合是非支配的。在这里，权重$\alpha(S)$就是这种非支配性的证明，而不是在(7)中作为逐期控制的要素。

\paragraph{支持性特征描述：命题1（第~页，第8-11行）。}
打印出的命题表明，集合$T$是非支配的，当且仅当存在某个值$v>0$，使得$T$是最优解。
\[
\max_{S\subset N}\{R(S)-vQ(S)\}.
\]
这里的标量$v$只是用于描述特征的辅助值，并非(7)中的额外决策。所展示的命题中使用了$S\subset N$，而定义1和该命题的证明则使用了$S\subseteq N$。这里的操作族完全遵循定义1，即$S\subseteq N$。所展示的特征描述仍然保留了其打印出的子集符号。

\paragraph{简化与边界情况：命题2（p.~11）。}
源代码中提到，如果$\Delta V_{t-1}(x)>0$是一个被支配的集合，那么在原始的无限制提供集合递归中，它永远不是最优的。而如果$\Delta V_{t-1}(x)=0$是一个非被支配的集合，那么总存在一个最优的非被支配集合。因此，(7)保留了非被支配的集合，而不是对其状态施加严格的积极边际价值限制。同时，也没有选择性地消除那些无法成为强非被支配集合的非被支配集合。

\paragraph{排序与系数：命题3以及前面的定义（7）（pp.~11--12）。}
让$m$表示非支配集合的数量，让$S_1,\ldots,S_m$用来枚举这些集合。对它们进行编号，使得……
\[
Q(S_1)\le Q(S_2)\le\cdots\le Q(S_m).
\]
命题 3 给出了相应的顺序。
\[
R(S_1)\le R(S_2)\le\cdots\le R(S_m).
\]
在(7)中的系数分别是：
\[
R_k=R(S_k),\qquad Q_k=Q(S_k),\qquad k=1,\ldots,m.
\]
这些是非严格的顺序。对于这种一般选择公式，不需要包含排序$S_1\subseteq\cdots\subseteq S_m$。

\paragraph{打印边界遗漏与限定。}
方程(2)如上所述，并未指定$V_0(0)$的值，尽管(7)中的$\Delta V_0(1)$指的是它。这里没有为该角点添加任何值。此外，命题2中的严格排除声明并未提及$\lambda>0$：在$\lambda=0$处，(7)中最大化表达式，以及在原始递归中，对于每个提供的集合来说都是相同的。这并不与命题相矛盾，因为命题中的严格排除声明假设了正的边际价值；对于打印模型，并没有施加额外的严格正性假设。

\paragraph{来源范围。}
本模型依据 Talluri 和 van Ryzin 的论文 \emph{基于一般离散选择消费者行为模型的收益管理}，日期为 2001 年 10 月 21 日，见 PDF 实际第 ~5 至 12 页。 方程（7）是所选的递归方式，它取代了那些不受限制的提供集合递归方式；这些递归方式的边际价值、概率、收益以及边界定义都继承自上述定义。 仅用于证明的线性规划问题、专门的选择模型、嵌套分配实现方式以及估计方法，均不属于本公式化的内容。

\end{document}
